\section{The minimal model revisited}\label{sec:min}
% src: L7 Sections 1 to 3, 5 and 6; answers to L7-E7, L7-E11, L7-E12, L7-O5 in the change lists of B4 and B5; dennis2026communicating

The research report of November 7, 2026 \citep{dennis2026communicating} studied a minimal version of the present framework: a linear-Gaussian world, agents that sense blocks of it, rank-limited messages in private linear encodings, attention and transmission chosen by expected free energy, and receivers that read messages as noisy observations of the world. It asked whether such agents condense into interior agents surrounded by a Markov blanket of agents, the picture behind hypothesis H1, and it obstructed that hypothesis without refuting it: message-fed blankets of agents exist as strict, locally stable orbits in four declared variants of the model, and none both passes the report's attribution rule and is reached from an uninformed start more often than under a null that randomizes partners. Its results are derivations, exhaustive enumerations and searches at declared parameter points, and mean-field computations; it simulated no sample path. This section states two results about the moment bound, one for a common rank-one code and one for noiseless grounding. The Supplement (\cref{sup:min}) places the earlier model inside the framework, sorts its results by how much they use the private encodings, restates the main ones with their scope, proves a third result about the moment bound, on its failure, and says which conclusions the framework's additions reopen and where they are now answered. Numbers quoted from the report are its own; \cref{chk:min} reproduces those on which this section relies.

The earlier model is the linear-Gaussian case of the agent model of \cref{def:fw-system} with an exogenous world, one round per world step, linear encodings and a single modality in every computed instance. The Supplement gives the dictionary between the two models and sorts the report's results into classes by how much of the encoding structure they use (\cref{sec:min-dictionary,def:min-classes}). It restates the report's message-fed bilayers, the four variants in which the report found them and the reasons none qualifies (\cref{sec:min-bilayers,def:min-bilayer}), and the report's attention structures, with a remark on why the reciprocal pairs of its replicator analysis need differences between frames (\cref{sec:min-attention,rem:min-reciprocity}).

\subsection{Double counting and the moment bound}\label{sec:min-moment}

Every length and time of the scale theory, and the definition of a collective, assume the moment bound (M): a stationary world-conditional law with finite second moments or, when attention changes from round to round, second moments bounded uniformly in time (\cref{asm:sca-standing,def:col-collective}); the report's version bounds the second moments of the decision-time belief means. The report found that its receivers' attention can violate it. This subsection states what restores the bound; the misreading behind that failure and the proof of when the failure starts are in the Supplement.

The misreading behind the failure is double counting: a level-1 receiver treats the error of a sender's posted predictive mean as independent of its own (Supplement, \cref{sup:min-double-counting-and-amplification}). The failure can start in a single round. When two agents that do not sense listen only to each other at level 1 with known rank-one codes, the reception step multiplies the homogeneous part of their disagreement by a matrix whose spectral radius exceeds one exactly under a criterion that proportional codes never meet (Supplement, \cref{lem:min-disagreement}, under the hypotheses stated there). Expected free energy selects a loop of mutual listening, in which both agents also sense, in the two-agent profile of the supplement (\cref{ex:sca-amplifying-profile}), whose frozen mean map has spectral radius $1.0613$ (\cref{chk:min-amplification}); this profile is the two-agent counterexample referred to below. It uses sender-specific codes, and one common code removes the amplification in two cases.

\begin{proposition}[One common rank-one code]\label{prop:min-common-code}
Let the effective codes be positive multiples of one nonzero row vector, $H_j=\vartheta_jT_c$ with $T_c\in\R^{1\times K}$ and $\vartheta_j>0$ (at identity encodings $\vartheta_j=1$), let the channel noise be $R_c=1$ (a sender-dependent noise variance can be absorbed into $\vartheta_j$), and let receivers know the codes and read at level 1. In one round let agent $i$ have decision-time covariance $P_i\succ0$ with precision $\Lambda_i=P_i^{-1}$, sensing information $\Gamma_i=(H^s_i)^\top(R^s_i)^{-1}H^s_i\succeq0$ (zero if it does not sense) and listening set $J_i$. Put $\alpha_j=\vartheta_j^2/(1+\vartheta_j^2T_cP_jT_c^\top)$ and $\Lambda^+_i=\Lambda_i+\Gamma_i+\sum_{j\in J_i}\alpha_jT_c^\top T_c$, and let $W^{\mathrm{rec}}$ be the homogeneous reception map
\begin{equation}\label{eq:min-reception}
e_i\mapsto(\Lambda^+_i)^{-1}\Bigl[\Lambda_ie_i+\sum_{j\in J_i}\alpha_jT_c^\top T_ce_j\Bigr],\qquad i\in V.
\end{equation}
\begin{enumerate}
\item[(a)] If listening is symmetric ($j\in J_i$ exactly when $i\in J_j$), then $W^{\mathrm{rec}}$ is self-adjoint for the inner product $\langle x,y\rangle_{\tilde M}=x^\top\tilde My$ with $\tilde M=\diag_i(\alpha_i\Lambda^+_i)$, it is diagonalizable with real spectrum in $(-1,1]$, and its eigenspace for the eigenvalue $1$ is
\[
\mathcal P=\bigl\{e:H^s_ie_i=0\text{ for every }i,\quad T_ce_i=T_ce_j\text{ whenever }j\in J_i\bigr\}.
\]
\item[(b)] If in addition the profile is frozen at a covariance fixed point and the world is isotropic, $F=fI_K$ with $0<f<1$, the error-transfer matrix $W=fW^{\mathrm{rec}}$, which is the matrix of \cref{def:fw-level0} with these level-1 gains, satisfies $\rho(W)\le f$ and $\Vert W^t\Vert_2\le K_{\tilde M}f^t$ for every $t\ge0$, with $K_{\tilde M}=(\lambda_{\max}(\tilde M)/\lambda_{\min}(\tilde M))^{1/2}$. Equality $\rho(W)=f$ holds exactly when $\mathcal P\neq\{0\}$, which is the case unless $\ker H^s_i\cap\ker T_c=\{0\}$ for every agent and every connected component of the listening graph contains an agent whose sensing map is injective. If channel noises are specific to each receiver, so that the local noise is independent across agents with blocks bounded by $q$, the world-conditional covariance satisfies
\[
\Vert\Sigma^c_{ij}\Vert\le\frac{K_{\tilde M}^2qf^{d(i,j)}}{1-f^2},
\]
with $d$ the hop distance on the listening graph.
\item[(c)] If no agent senses, then $\rho(W^{\mathrm{rec}})=1$ for every listening relation, symmetric or not, and every choice of the scales $\vartheta_j$. In the world $F=fI_K$, at a covariance fixed point of a frozen profile $\rho(W)=f=\rho(F)$, and along a periodic covariance orbit of a profile of period $p$ the map over one period has spectral radius $f^p$.
\end{enumerate}
\end{proposition}
% src: L7-P2, L7-P3 (a), (b), (b'), (c); B3-P1 row (ii) and B3-P2(a) for the spectral form; the periodic case of (c) and the characterization of P = {0} are added here
Proof in \cref{app:min}; the spectral statements were also checked numerically (\cref{chk:min-common-code}). The subspace $\mathcal P$ holds the error configurations that no agent's sensing sees and on which neighbors' coded messages agree, so no message can correct them: agreement in expression with disagreement in private opinion, as in discourse-sheaf models of opinion dynamics \citep{hansen2021opinion}, here an unobservable error mode. By part (b), with one code a community relaxes faster than the world only if every member senses all directions but at most one, which the code must read, and some member of each listening component senses every direction. Part (c) proves, for one common rank-one code in an isotropic world, the report's conjecture that a frozen profile in which no agent senses has a mean map of spectral radius below one, with the value $f$: a pure-listening community relaxes only at the world's own rate, whether its attention is frozen or switches periodically. The clustering bound of part (b) holds where the row-sum contraction of the clustering theorem fails in most instances (377 of 400 in \cref{chk:min-common-code}), which is why the scale theory needs its spectral form (\cref{prop:sca-spectral}).

Parts (b) and (c) prove two cases of \cref{conj:sca-common-code}; the cases, the computations that support the conjecture and the return of sender-specific codes with several modalities are discussed in the Supplement (\cref{sup:min-double-counting-and-amplification}). Noisy sensing does not restore (M), since both agents of the two-agent counterexample sense every round. Noise-free observation of the world does, above a rate.

\begin{proposition}[Noiseless grounding restores the moment bound]\label{prop:min-grounding-moment}
Let $Q\succ0$, let receivers know the senders' effective codes $H_j$, of any rank, and read at level 1 with assigned noise $\hat R_{ij}=R_c+H_jP_jH_j^\top$, $R_c\succ0$, the channel noise being independent across senders and of the sensing noise. Let attention depend on the past in any way, with at most $\bar J$ receivers listening to any one sender in a round; let $\eta_0\sim\Ncal(0,P_0)$, let the initial errors have finite second moments and let $Q\preceq P_i(0)\preceq P_0$ for every $i$. Suppose that in each round each agent, independently of everything else and with probability $p_g$, observes $\eta_t$ without noise, so that its posterior error is zero. Put
\begin{gather*}
\varpi_Q=1-\frac{1}{\lambda_{\max}(Q^{-1}P_0)},\qquad\mathrm{snr}_j=\lambda_{\max}\bigl(R_c^{-1/2}H_jP_0H_j^\top R_c^{-1/2}\bigr),\\
\varpi_H=\max_j\frac{\mathrm{snr}_j}{1+\mathrm{snr}_j},\qquad\varpi_g=(1-p_g)\varpi_Q(1+\varpi_H\bar J).
\end{gather*}
Then $\mathcal V(t)=\sum_ie_i(t)^\top\Lambda_i(t)e_i(t)$ satisfies $\E\mathcal V(t+1)\le\varpi_g\E\mathcal V(t)+N\tr(P_0)/\lambda_{\min}(Q)$, and the part of the error recursion that does not involve new noise contracts $\mathcal V$ in mean by the factor $\varpi_g$. If $\varpi_g<1$, the second moments of the errors and of the decision-time belief means are bounded uniformly in time, which is (M) in its uniform form.
\end{proposition}
% src: L7-P4 (the report's projection, header and prediction-contraction lemmas are reproved in the appendix); B3-P1 row (vi)
Proof in \cref{app:min}. In the two-agent counterexample $\varpi_Q=0.81$, $\mathrm{snr}_j=9/0.19$, $\varpi_H=0.97933$ and $\bar J=1$, so grounding each agent with probability above $0.3763$ per round restores the bound (\cref{chk:min-grounding-moment}). At $p_g=0$ the proposition is a small-gain condition like the report's, except that the report bounds the largest weighted error and so carries the capacity $\kappa$, while the sum over agents carries the number $\bar J$ of listeners of one sender, and only the sum lets independent grounding help. The conclusion applies when attention changes from round to round; it gives neither pathwise contraction nor covariance clustering. The report's grounding rounds, which show a public referent unrelated to $\eta$ so that translations can be identified, carry no information about $\eta$ and do not restore (M); joint-attention grounding, a noisy observation of the attended part of $\eta$, lies between the two cases and is open.

The Supplement restates the report's group-agency results and the two formalizations of the preference to predict and be predicted, between which the framework leaves the choice open (\cref{sec:min-agency}). It also lists the conclusions of the report that grounding, records, performative acts and multiplexity may change and where they now stand: three are answered within a stated scope, eleven in part, and three remain open (\cref{sec:min-reexamination}). The calculation the report expected to change its verdict becomes a question about collectives, which is an open problem (\cref{open:min-exact-fusion}).
